% Modul Belajar 11, dihasilkan oleh tools/build.py dari tools/konten/p11.py
\documentclass[a4paper]{article}
\usepackage{modulITERA}
\begin{document}
\Sampul{11}{Searching}{Mencari data di array dengan linear search dan binary search, dan membandingkan banyak langkahnya}{CPMK0607 (menerapkan) dan CPMK0608 (menjelaskan) pada konsep dasar algoritma dan sistem komputer}{sekitar 3 jam belajar mandiri, ditambah 150 menit di kelas}{\texttt{02 Hands-on/P11 - Searching - Hands-on/}}
\Bagian{Modul 11}{Peta belajar}
Ikuti urutan ini dari atas ke bawah. Setiap bagian memakai apa yang sudah dibahas di bagian sebelumnya.
\par\vspace{4pt}\noindent{\fontsize{9.5pt}{13pt}\selectfont
\begin{tabular}{@{}>{\raggedright\arraybackslash}p{0.140\textwidth}>{\raggedright\arraybackslash}p{0.840\textwidth}@{}}
\kepala{Urutan} & \kepala{Bagian}\\
1 & Tentang modul ini\\\garis
2 & Masalah pembuka: Mencari nama di daftar wisuda\\\garis
3 & Langkah 1: Linear search\\\garis
4 & Langkah 2: Mencari semua kemunculan\\\garis
5 & Langkah 3: Ide binary search\\\garis
6 & Langkah 4: Binary search di C++\\\garis
7 & Langkah 5: Berapa langkahnya?\\\garis
8 & Pesan error yang sering muncul\\\garis
9 & Latihan mandiri\\\garis
10 & Rangkuman\\\garis
11 & Cek pemahaman\\\garis
12 & Exit Ticket dan tugas\\\garis
13 & Bacaan lanjut dan persiapan minggu depan\\
\garisdasar
\end{tabular}}\par\vspace{6pt}
\Bagian{Pembuka}{Tentang modul ini}
Modul ini mendampingi Pertemuan 11. Sejauh ini kalian sudah menyimpan banyak data di array. Pertanyaan berikutnya: bagaimana menemukan satu data di antaranya dengan cepat? Minggu ini kalian membandingkan dua cara mencari, dan melihat mengapa data yang terurut bisa dicari jauh lebih cepat.

Isinya sama dengan slide di kelas, tetapi ditulis lebih pelan dan lebih lengkap, supaya kalian bisa mengulang sendiri di rumah tanpa harus menebak maksud slide.

\Sub{Setelah menyelesaikan modul ini, kalian bisa}
\par\vspace{4pt}\noindent{\fontsize{9.5pt}{13pt}\selectfont
\begin{tabular}{@{}>{\raggedright\arraybackslash}p{0.070\textwidth}>{\raggedright\arraybackslash}p{0.680\textwidth}>{\raggedright\arraybackslash}p{0.230\textwidth}@{}}
\kepala{No} & \kepala{Kemampuan} & \kepala{Dibahas di}\\
1 & Menulis program yang mencari data pada array dengan linear search dan binary search, termasuk mencari semua kemunculan dan menangani data yang tidak ditemukan (Sub-CPMK 11.1, CPMK0607) & Langkah 1 sampai 5\\\garis
2 & Menelusuri langkah binary search dengan trace table low, high, dan mid, serta menjelaskan perbedaan banyak perbandingan kedua algoritma (Sub-CPMK 11.2, CPMK0608) & kotak Tebak dulu, Latihan 3\\
\garisdasar
\end{tabular}}\par\vspace{6pt}
Karena setiap instrumen penilaian dibagi rata antara Bagian A (menulis kode) dan Bagian B (menelusuri dan menjelaskan), yang dinilai bukan hanya program kalian jalan, tetapi juga kalian bisa menjelaskan mengapa program itu jalan.

\Sub{Cara memakai modul ini}
Setiap langkah punya pola yang sama: penjelasan konsep, kadang contoh dari kehidupan sehari-hari, lalu kode yang bisa langsung kalian ketik. Sesudah kode ada kotak \textbf{Tebak dulu}. Tulis tebakan kalian di kertas sebelum membaca \textbf{Jawaban} di bawahnya. Kebiasaan menebak lalu membuktikan inilah yang membuat konsep menempel, bukan sekadar membaca.

\begin{Penting}[Ketik, jangan salin]
Ketik ulang setiap contoh di GDB Online atau editor kalian, jangan disalin-tempel. Saat mengetik, kalian akan membuat salah ketik, membaca pesan error, dan memperbaikinya. Itu bagian dari belajar.
\end{Penting}
\Sub{Yang perlu disiapkan}
Peramban dengan akses ke onlinegdb.com, atau g++ di komputer kalian sendiri. Kertas dan alat tulis untuk menebak keluaran dan membuat trace table. Folder hands-on pertemuan ini dari repo mata kuliah.

\Bagian{Pembuka}{Masalah pembuka: Mencari nama di daftar wisuda}
Panitia wisuda memegang daftar 1.200 nama yang sudah urut abjad. Seorang orang tua bertanya, apakah anaknya ada di daftar?

Membaca dari nama pertama sampai ketemu bisa butuh ratusan kali melihat. Padahal saat mencari kata di kamus, kita tidak membaca dari halaman pertama.

\begin{MKeluaran}[title={Tampilan yang diinginkan}]
Cari: Siti Rahma
Ditemukan di nomor 1023
Linear search : 1023 perbandingan
Binary search : 10 perbandingan
\end{MKeluaran}
\begin{Tebak}[Coba pikirkan]
Bagaimana kalian mencari kata di kamus? Mengapa cara itu tidak bisa dipakai pada daftar yang belum urut?
\end{Tebak}
\begin{Jawaban}[Cara memecahnya]
Bandingkan satu per satu: Cara paling sederhana: periksa setiap elemen dari depan. Manfaatkan urutan: Bila data urut, buang separuh daftar di setiap langkah. Terjemahkan: Linear search dengan \texttt{for}, binary search dengan \texttt{while} dan tiga indeks.
\end{Jawaban}
\Bagian{Langkah 1}{Linear search}
Linear search memeriksa elemen satu per satu dari depan. Begitu cocok, posisinya dicatat dan pencarian berhenti dengan \texttt{break}. Bila sampai akhir tidak ada yang cocok, variabel posisi tetap -1.

Linear search bekerja pada data apa pun, terurut atau tidak. Kelemahannya, pada kasus terburuk semua elemen diperiksa.

\begin{Contoh}[Contoh sehari-hari]
Mencari kunci motor di tas yang berantakan: kalian mengeluarkan barang satu per satu sampai kuncinya ketemu, atau sampai tas kosong.
\end{Contoh}
\begin{MKode}{linear.cpp}
int a[8] = {4, 9, 2, 7, 5, 9, 1, 8};
int x = 7, posisi = -1;
for (int i = 0; i < 8; i++) {
    if (a[i] == x) {
        posisi = i;
        break;
    }
}
if (posisi == -1) cout << "tidak ada" << endl;
else cout << "indeks " << posisi << endl;
\end{MKode}
\begin{Tebak}[]
Sebelum menjalankan program, tulis keluarannya di kertas, karakter demi karakter.
\end{Tebak}
\begin{Jawaban}[]
\begin{MKeluaran}[]
indeks 3
\end{MKeluaran}
Nilai awal -1 menandai belum ditemukan, karena -1 tidak pernah menjadi indeks yang sah.
\end{Jawaban}
\Bagian{Langkah 2}{Mencari semua kemunculan}
Kadang yang ditanyakan bukan di mana pertama kali, tetapi di mana saja dan berapa kali. Untuk itu \texttt{break} dihilangkan dan sebuah pencacah ditambahkan.

\begin{MKode}{semua.cpp}
int a[8] = {4, 9, 2, 7, 5, 9, 1, 8};
int x = 9, banyak = 0;
for (int i = 0; i < 8; i++) {
    if (a[i] == x) {
        cout << "indeks " << i << endl;
        banyak++;
    }
}
cout << "muncul " << banyak << " kali" << endl;
\end{MKode}
\begin{Tebak}[]
Sebelum menjalankan program, tulis keluarannya di kertas, karakter demi karakter.
\end{Tebak}
\begin{Jawaban}[]
\begin{MKeluaran}[]
indeks 1
indeks 5
muncul 2 kali
\end{MKeluaran}
Untuk semua kemunculan, perulangan tidak dihentikan; pencacah mencatat banyaknya.
\end{Jawaban}
\Bagian{Langkah 3}{Ide binary search}
Bila data sudah terurut, kita tidak perlu memeriksa dari depan. Periksa elemen tengah. Bila yang dicari lebih kecil, buang separuh kanan; bila lebih besar, buang separuh kiri. Setiap langkah membuang separuh sisa data.

Tiga indeks menjaga batas: \texttt{low} dan \texttt{high} mengapit bagian yang masih mungkin, \texttt{mid} adalah tengahnya. Trace table berikut menunjukkan pencarian 21.

\begin{Contoh}[Contoh sehari-hari]
Tebak angka 1 sampai 100. Penebak yang cerdik selalu menebak tengah, lalu membuang separuh kemungkinan setiap kali diberi tahu terlalu besar atau terlalu kecil.
\end{Contoh}
\par\vspace{4pt}\noindent{\fontsize{9.5pt}{13pt}\selectfont
\begin{tabular}{@{}>{\raggedright\arraybackslash}p{0.131\textwidth}>{\raggedright\arraybackslash}p{0.098\textwidth}>{\raggedright\arraybackslash}p{0.098\textwidth}>{\raggedright\arraybackslash}p{0.098\textwidth}>{\raggedright\arraybackslash}p{0.131\textwidth}>{\raggedright\arraybackslash}p{0.425\textwidth}@{}}
\kepala{Langkah} & \kepala{low} & \kepala{high} & \kepala{mid} & \kepala{a[mid]} & \kepala{Keputusan}\\
1 & 0 & 9 & 4 & 34 & 34 > 21, cari di kiri: high = 3\\\garis
2 & 0 & 3 & 1 & 8 & 8 < 21, cari di kanan: low = 2\\\garis
3 & 2 & 3 & 2 & 15 & 15 < 21, low = 3\\\garis
4 & 3 & 3 & 3 & 21 & ditemukan di indeks 3\\
\garisdasar
\end{tabular}}\par\vspace{6pt}
\begin{Penting}[]
Binary search hanya benar pada data yang sudah terurut.
\end{Penting}
\Bagian{Langkah 4}{Binary search di C++}
Terjemahan langsung dari trace table. Perhatikan tiga detail yang sering salah: syarat \texttt{low <= high}, pergeseran \texttt{mid + 1} dan \texttt{mid - 1} (bukan \texttt{mid}), dan arah pergeseran yang mengikuti hasil perbandingan.

\begin{MKode}{biner.cpp}
int a[10] = {3, 8, 15, 21, 34, 42, 57, 66, 79, 90};
int x = 21, low = 0, high = 9, posisi = -1;
while (low <= high) {
    int mid = (low + high) / 2;
    if (a[mid] == x) { posisi = mid; break; }
    else if (a[mid] < x) low = mid + 1;
    else high = mid - 1;
}
cout << "indeks " << posisi << endl;
\end{MKode}
\begin{Tebak}[]
Sebelum menjalankan program, tulis keluarannya di kertas, karakter demi karakter.
\end{Tebak}
\begin{Jawaban}[]
\begin{MKeluaran}[]
indeks 3
\end{MKeluaran}
Syaratnya \texttt{low <= high}. Dengan \texttt{<}, elemen terakhir yang tersisa tidak sempat diperiksa.
\end{Jawaban}
\Bagian{Langkah 5}{Berapa langkahnya?}
Pada kasus terburuk, linear search memeriksa semua n elemen. Binary search hanya butuh sekitar log₂ n langkah, karena data terbagi dua di setiap langkah. Selisihnya makin besar ketika datanya makin banyak.

\par\vspace{4pt}\noindent{\fontsize{9.5pt}{13pt}\selectfont
\begin{tabular}{@{}>{\raggedright\arraybackslash}p{0.283\textwidth}>{\raggedright\arraybackslash}p{0.348\textwidth}>{\raggedright\arraybackslash}p{0.348\textwidth}@{}}
\kepala{Banyak data} & \kepala{Linear, kasus terburuk} & \kepala{Binary, kasus terburuk}\\
10 & 10 & 4\\\garis
100 & 100 & 7\\\garis
1.000 & 1.000 & 10\\\garis
1.000.000 & 1.000.000 & 20\\
\garisdasar
\end{tabular}}\par\vspace{6pt}
\begin{Penting}[]
Binary search jauh lebih cepat, tetapi butuh data terurut. Mengurutkan data adalah materi minggu depan.
\end{Penting}
\Bagian{Waspada}{Pesan error yang sering muncul}
Semua pesan di tabel ini dihasilkan g++ dengan opsi \texttt{-Wall}, sama seperti yang dipakai \texttt{cek.sh}.
\par\vspace{4pt}\noindent{\fontsize{9.5pt}{13pt}\selectfont
\begin{tabular}{@{}>{\raggedright\arraybackslash}p{0.240\textwidth}>{\raggedright\arraybackslash}p{0.270\textwidth}>{\raggedright\arraybackslash}p{0.470\textwidth}@{}}
\kepala{Kesalahan} & \kepala{Contoh} & \kepala{Pesan pertama dari g++}\\
Membandingkan array dengan == & \texttt{if (a == x)} & \texttt{error: ISO C++ forbids comparison between pointer and integer}\\\garis
Variabel di luar jangkauan & \texttt{mid dipakai di luar while} & \texttt{error: 'mid' was not declared in this scope}\\\garis
Kutip tunggal untuk string & \texttt{nama[i] == 'Rani'} & \texttt{error: no match for 'operator=='}\\\garis
Posisi tidak dipakai & \texttt{int posisi = -1; tak dipakai} & \texttt{warning: unused variable 'posisi'}\\\garis
Kurang titik koma di dalam if & \texttt{posisi = mid break;} & \texttt{error: expected ';' before 'break'}\\
\garisdasar
\end{tabular}}\par\vspace{6pt}
Kesalahan binary search yang paling sering, yaitu syarat \texttt{low < high} dan arah pergeseran yang terbalik, tidak memunculkan pesan apa pun. Program tetap jalan tetapi jawabannya salah. Trace table adalah satu-satunya cara cepat menemukannya.

\Bagian{Latihan}{Latihan mandiri}
Latihan ini sama dengan hands-on di kelas. Semua file ada di \texttt{02 Hands-on/P11 - Searching - Hands-on/latihan/}. Kerjakan berurutan, lalu cocokkan dengan \texttt{expected/} atau jalankan \texttt{bash cek.sh}.
\par\vspace{4pt}\noindent{\fontsize{9.5pt}{13pt}\selectfont
\begin{tabular}{@{}>{\raggedright\arraybackslash}p{0.070\textwidth}>{\raggedright\arraybackslash}p{0.360\textwidth}>{\raggedright\arraybackslash}p{0.150\textwidth}>{\raggedright\arraybackslash}p{0.400\textwidth}@{}}
\kepala{No} & \kepala{File} & \kepala{Tingkat} & \kepala{Yang dilatih}\\
1 & \texttt{handson1-cari.cpp} & Dasar & linear search, hasil -1, hitung perbandingan\\\garis
2 & \texttt{handson2-nama.cpp} & Menengah & mencari semua kemunculan string\\\garis
3 & \texttt{handson3-biner.cpp} & Tantangan & trace low-high-mid, perbaiki tiga kesalahan\\\garis
+ & \texttt{bonus-banding.cpp} & Pengayaan & menghitung perbandingan dua algoritma\\
\garisdasar
\end{tabular}}\par\vspace{6pt}
\Sub{Latihan 1: handson1-cari.cpp}
Linear search dengan kasus ditemukan dan tidak ditemukan, serta banyak perbandingan.

\Sub{Latihan 2: handson2-nama.cpp}
Cari semua posisi sebuah nama di daftar hadir dan banyak kemunculannya.

\Sub{Latihan 3: handson3-biner.cpp}
Binary search dengan tiga kesalahan. Buat trace table untuk dua pencarian, perbaiki, dan jelaskan.

\Sub{Bonus: bonus-banding.cpp}
Bandingkan banyak perbandingan linear dan binary search untuk empat nilai.

\Sub{Latihan menelusuri}
\begin{MKode}{tebak.cpp}
int a[7] = {2, 5, 8, 12, 16, 23, 38};
int x = 23, low = 0, high = 6, langkah = 0;
while (low <= high) {
    int mid = (low + high) / 2;
    langkah++;
    if (a[mid] == x) break;
    if (a[mid] < x) low = mid + 1;
    else high = mid - 1;
}
cout << langkah << " " << low << " " << high << endl;
\end{MKode}
\begin{Tebak}[]
Tulis keluarannya di kertas tanpa menjalankan program. Buat trace table bila perlu.
\end{Tebak}
\begin{Jawaban}[]
\begin{MKeluaran}[]
2 4 6
\end{MKeluaran}
mid = 3 (12 < 23, low = 4), lalu mid = 5 (23, ketemu). Dua langkah; saat berhenti low = 4 dan high = 6.
\end{Jawaban}
\Bagian{Penutup}{Rangkuman}
\par\vspace{4pt}\noindent{\fontsize{9.5pt}{13pt}\selectfont
\begin{tabular}{@{}>{\raggedright\arraybackslash}p{0.320\textwidth}>{\raggedright\arraybackslash}p{0.660\textwidth}@{}}
\kepala{Konsep} & \kepala{Yang perlu diingat}\\
Linear search & Nilai awal -1 menandai belum ditemukan, karena -1 tidak pernah menjadi indeks yang sah.\\\garis
Mencari semua kemunculan & Untuk semua kemunculan, perulangan tidak dihentikan; pencacah mencatat banyaknya.\\\garis
Ide binary search & Binary search hanya benar pada data yang sudah terurut.\\\garis
Binary search di C++ & Syaratnya \texttt{low <= high}.\\\garis
Berapa langkahnya? & Binary search jauh lebih cepat, tetapi butuh data terurut.\\
\garisdasar
\end{tabular}}\par\vspace{6pt}
\Bagian{Penutup}{Cek pemahaman}
Jawab tanpa menjalankan program, lalu periksa dengan menjalankannya.
\begin{enumerate}
\item Mengapa posisi diawali -1?
\item Berapa perbandingan linear search untuk mencari elemen terakhir dari 50 data?
\item Apa prasyarat binary search?
\item Mengapa syaratnya \texttt{low <= high}, bukan \texttt{low < high}?
\item Kira-kira berapa langkah binary search untuk 1.000 data?
\end{enumerate}
\begin{Jawaban}[]
\begin{enumerate}[leftmargin=16pt]
\item -1 tidak pernah menjadi indeks yang sah, jadi menandai belum ditemukan.
\item 50.
\item Data harus sudah terurut.
\item Saat low sama dengan high masih ada satu elemen yang harus diperiksa.
\item Sekitar 10 langkah.
\end{enumerate}
\end{Jawaban}
\Bagian{Penilaian}{Exit Ticket 11}
Dikerjakan di 25 menit terakhir kelas, sendiri, tanpa AI, internet, atau diskusi. Exit Ticket 11 adalah satu dari 14 Exit Ticket; rata-ratanya menjadi komponen berbobot 25\%.
\par\vspace{4pt}\noindent{\fontsize{9.5pt}{13pt}\selectfont
\begin{tabular}{@{}>{\raggedright\arraybackslash}p{0.080\textwidth}>{\raggedright\arraybackslash}p{0.600\textwidth}>{\raggedright\arraybackslash}p{0.100\textwidth}>{\raggedright\arraybackslash}p{0.200\textwidth}@{}}
\kepala{Soal} & \kepala{Isi} & \kepala{Poin} & \kepala{CPMK}\\
A1 & Tulis linear search yang menampilkan semua indeks kemunculan sebuah nilai & 25 & CPMK0607\\\garis
A2 & Tulis binary search pada array terurut yang menampilkan indeks atau tidak ditemukan & 25 & CPMK0607\\\garis
B1 & Isi trace table low, high, mid untuk sebuah pencarian & 20 & CPMK0608\\\garis
B2 & Jelaskan mengapa binary search gagal pada data yang belum terurut & 20 & CPMK0608\\\garis
B3 & Bandingkan banyak perbandingan dua algoritma untuk data yang diberikan & 10 & CPMK0608\\
\garisdasar
\end{tabular}}\par\vspace{6pt}
\Bagian{Penilaian}{Tugas 11: Direktori Kontak}
Kumpulkan sebelum Pertemuan 12 lewat kanal pengumpulan yang diumumkan dosen.

\Sub{Bagian A: program (50 poin)}
Buat program direktori kontak dengan menu berikut. Data dimasukkan dalam urutan abjad nama; program menolak nama yang tidak urut.
\begin{enumerate}
\item Tambah kontak, maksimal 10 (nama satu kata dan nomor telepon), dengan validasi urutan abjad
\item Tampilkan semua kontak
\item Cari nama dengan linear search; tampilkan nomor dan banyak perbandingan
\item Cari nama dengan binary search; tampilkan nomor dan banyak perbandingan
\item Keluar
\end{enumerate}
\begin{Contoh}[Bonus, tidak wajib]
Tambahkan pencarian awalan nama, misalnya semua kontak yang diawali "Ra".
\end{Contoh}
\Sub{Bagian B: penjelasan (50 poin)}
Komentar maksimal 150 kata: trace table binary search untuk satu nama yang ada dan satu yang tidak ada, serta perbandingan banyak langkah linear dan binary search pada data kalian. Jelaskan mengapa data harus urut abjad.

\Sub{Rubrik Tugas 11}
\par\vspace{4pt}\noindent{\fontsize{8.5pt}{11.5pt}\selectfont
\begin{tabular}{@{}>{\raggedright\arraybackslash}p{0.190\textwidth}>{\raggedright\arraybackslash}p{0.070\textwidth}>{\raggedright\arraybackslash}p{0.200\textwidth}>{\raggedright\arraybackslash}p{0.180\textwidth}>{\raggedright\arraybackslash}p{0.170\textwidth}>{\raggedright\arraybackslash}p{0.170\textwidth}@{}}
\kepala{Kriteria} & \kepala{Poin} & \kepala{Sangat baik 85–100} & \kepala{Baik 70–84} & \kepala{Cukup 55–69} & \kepala{Kurang <55}\\
A1 Program berjalan & 20 & tanpa error dan peringatan & ada peringatan & perlu satu perbaikan & tidak terkompilasi\\\garis
A2 Kelengkapan menu & 15 & lima menu dan validasi urutan & satu menu kurang & dua kurang & tanpa binary search\\\garis
A3 Ketepatan pencarian & 15 & kedua pencarian tepat, termasuk tidak ditemukan & satu kasus salah & dua salah & pencarian salah\\\garis
B1 Trace table dan perbandingan & 30 & dua trace tepat dan analisis jelas & satu trace keliru & tanpa analisis & tidak ada atau disalin\\\garis
B2 Cerita error dan pengujian & 20 & pesan, sebab, perbaikan, uji & pesan tidak dikutip & hanya perbaikan & tidak ada\\
\garisdasar
\end{tabular}}\par\vspace{6pt}
Nilai kriteria = poin × skor level. Contoh: A1 di level Baik dengan skor 80 memberi 20 × 0,80 = 16 poin.

\Bagian{Penutup}{Bacaan lanjut dan persiapan minggu depan}
\begin{enumerate}
\item Deitel, P. dan Deitel, H. C++ How to Program, edisi ke-10, Bab 8 (searching).
\item learncpp.com, Bab 18 (algoritma pada array).
\item Visualgo.net, modul Searching dan Sorting.
\item Slide \texttt{01 Slide/P11 Searching.pdf} dan hands-on di \texttt{02 Hands-on/P11 - Searching - Hands-on/}.
\end{enumerate}
\begin{Penting}[Minggu depan]
Pertemuan 12 membahas sorting: bubble sort dan selection sort, supaya data bisa diurutkan sebelum dicari dengan binary search.
\end{Penting}
\end{document}
